\chapter{Individu yang Unik}\label{ch:g9:unique-individuals}

Tujuh miliar orang, dan tidak ada dua yang serupa --- tidak
di dalam sebuah keluarga, tidak di sebuah halaman sekolah,
tidak dalam seluruh sejarahnya, dengan kembar identik
setengah dikecualikan. Tiga bab kini sudah merakit seluruh
permesinan kenyataan itu. Bab penutup yang pendek ini
menjalankan hujahnya sekali, dari ujung ke ujung, lalu
menyimpan akibatnya: bagi keragaman spesiesnya, bagi cara
kita membaca tajuk berita ``genetik'', dan bagi apa yang
dimaksud serta tidak dimaksud oleh keunikan.

\section{Teorema keunikannya}

\begin{proposition}[Mengapa tidak ada dua orang yang serupa]\label{prop:g9:unique-individuals:theorem}
Rakitlah ketiga mekanismenya:
\begin{enumerate}
  \item \emph{pengocokan pada pembagian duanya}: \omterm{def:g8:sexual-reproduction:gametes}{gamet} tiap
        orang tua mengangkut satu mitra tiap pasang,
        dibagikan secara bebas --- $2^{23}$ setengah
        tumpukan tiap orang tua
        (\cref{rem:g9:chromosomes:whyhalve}), dan sebenarnya
        jauh lebih banyak, sebab mitra sepasangnya juga
        \emph{bertukar ruas} selama pembelahan pembagi
        duanya, memotong tumpukannya di bawah tingkat
        \omterm{def:g9:chromosomes:chromosomes}{kromosom} utuhnya;
  \item \emph{lotere \omterm{def:g5:flower-to-fruit:steps}{pembuahannya}}: sperma yang mana dari
        jutaannya (\cref{ex:g8:reproductive-systems:sperm})
        yang menemui \omterm{def:g2:animals-grow-up:born}{telur} bulan yang mana mengalikan kedua
        ruangnya;
  \item \emph{tetesan \omterm{prop:g9:genes-and-dna:explains}{mutasinya}}: setiap angkatan
        menambahkan beberapa ejaan segarnya sendiri
        (\cref{prop:g9:genes-and-dna:explains}).
\end{enumerate}
Ruang gabungannya mengerdilkan jumlah manusia yang pernah
hidup: setiap \omterm{def:g5:flower-to-fruit:steps}{pembuahan} membagikan sebuah tumpukan yang
belum pernah dibagikan, dan tidak akan pernah lagi.
\end{proposition}
\admitted

\begin{example}[Kembarnya, dan batas pengecualiannya]\label{ex:g9:unique-individuals:twins}
Kembar identik berbagi satu pembagian
(\cref{exo:g8:from-egg-to-birth:10}) --- dan bahkan mereka
pun memencar: kehidupan yang terpisah menuliskan \omterm{def:g9:heredity-and-traits:traits}{sifat}
perolehan yang berbeda
(\cref{prop:g9:heredity-and-traits:sorting}), riwayat
\omterm{def:g8:nervous-communication:synapse}{sinapsis} yang berbeda
(\cref{ex:g8:nervous-communication:juggler}), bahkan
beberapa \omterm{prop:g9:genes-and-dna:explains}{mutasi} lanjut yang berbeda di dalam \omterm{def:g5:first-look-at-cells:cell}{sel} tubuh
masing-masing. Satu-satunya retakan pada keunikannya menutup
dirinya sendiri: sesudah pembagiannya, kehidupanlah yang
membedakan.
\end{example}

\begin{remark}[Unik, dan sebagian besar serupa]\label{rem:g9:unique-individuals:alike}
Peganglah kedua kebenarannya: teks \omterm{def:g9:genes-and-dna:dna}{DNA} dua manusia mana pun
sepakat lebih dari $99$ persen --- pembangunan bersama satu
\omterm{def:g5:biodiversity:species}{spesies} pada \cref{prop:g5:first-look-at-cells:principle}
--- dan ketidaksepakatannya, yang tersebar menembus tiga
miliar huruf, tetap cukup untuk keunikan yang ketat.
``Semuanya berbeda'' dan ``semuanya sekerabat'' adalah
kenyataan yang sama pada dua perbesaran
(\cref{met:g9:genes-and-dna:levels}); setiap pemakaian
kebenaran yang pertama terhadap yang kedua adalah \omterm{def:g1:living-or-not:living}{biologi}
yang buruk sebelum ia menjadi etika yang buruk.
\end{remark}

\section{Keunikan yang bekerja}

\begin{example}[Pengenalan lewat DNA]\label{ex:g9:unique-individuals:forensics}
Keunikannya dapat diperiksa, dan pengadilan memeriksanya:
cukup banyak kedudukan yang berubah-ubah pada seseorang,
dibaca dari sejumput \omterm{def:g5:first-look-at-cells:cell}{selnya}, cocok dengan satu individu
(kembar identik dikesampingkan) dari seluruh umat manusia.
Perbandingan yang sama, dijalankan pada perbesaran
keluarganya, memecahkan soal keayahan --- tiap kedudukan
berubah-ubah pada anaknya harus terbagikan dari tumpukan
kedua orang tuanya. Logika album pada
\cref{ch:g9:heredity-and-traits} sudah menjadi sebuah alat.
\end{example}

\begin{example}[Transfusi dan pencangkokan]\label{ex:g9:unique-individuals:medicine}
Kedokteran berunding dengan keunikannya setiap hari.
Golongan \omterm{prop:g4:journey-of-food:absorb}{darah} (\cref{ex:g9:genes-and-dna:bloodgroups})
adalah kasus yang mudah --- empat kelas yang lebar, dapat
dicocokkan dari daftarnya. Pencangkokan organ menghadapi
kasus yang sukar: sistem kekebalannya
(\cref{ch:g9:immune-defenses} menjelaskan permesinannya)
membaca penanda-diri yang lebih halus butirannya, cukup unik
sehingga pendonornya dicari dari keluarganya lebih dahulu
--- makin dekat \omterm{prop:g6:classification-kinship:kinship}{kekerabatannya}, makin dekat pembagiannya.
\end{example}

\begin{proposition}[Keragaman adalah cadangan spesiesnya, tersimpan]\label{prop:g9:unique-individuals:reserve}
\Cref{rem:g8:sexual-reproduction:reserve} menjanjikan maksud
pengocokannya; mekanismenya kini menguangkannya. Sebuah
\omterm{def:g5:biodiversity:species}{spesies} yang terdiri atas individu yang unik memiliki,
tersebar di antara mereka, \omterm{def:g9:genes-and-dna:gene}{alel} untuk setiap keadaan ---
ketahanan yang diangkut sebagian terhadap penyakit ini,
kelapangan yang diangkut sebagian lain terhadap tekanan itu
(kebun \omterm{prop:g3:flowers-fruits-seeds:transform}{buah} pada \cref{ex:g8:sexual-reproduction:bets}, pada
skala \omterm{def:g5:biodiversity:species}{spesies}). Ketika keadaannya berbalik, penahan pada
\cref{prop:g8:reproduction-and-environment:limits} jatuh
tidak merata menembus keragaman itu --- dan apa yang
dilakukan ketidakmerataan itu, angkatan demi angkatan,
justru itulah pokok yang dibuka dua bab berikutnya:
berubahnya \omterm{def:g5:biodiversity:species}{spesies}.
\end{proposition}
\admitted

\begin{method}[Mengaudit sebuah klaim ``genetik'']\label{met:g9:unique-individuals:claims}
Untuk tajuk berita apa pun tentang \omterm{def:g9:genes-and-dna:gene}{gen} dan manusia:
\begin{enumerate}
  \item periksa tingkatnya
        (\cref{met:g9:genes-and-dna:levels}): apakah
        klaimnya tentang sebuah \omterm{def:g9:genes-and-dna:gene}{gen}, sebuah tumpukan, atau
        sebuah pola keluarga?
  \item periksa tenunannya
        (\cref{prop:g9:heredity-and-traits:sorting}): apakah
        ia menghormati tenunan
        rentang-turunan-yang-dijalani, atau berpura-pura
        satu \omterm{def:g9:genes-and-dna:gene}{gen} sama dengan satu takdir?
  \item periksa keragamannya: apakah ia memperlakukan
        sebaran sebuah \omterm{def:g8:reproduction-and-environment:population}{populasi} sebagai cadangan sebagaimana
        adanya, atau sebagai penyimpangan dari sebuah
        tumpukan ``normal'' yang tidak ada?
  \item periksa kembarnya: akankah ia bertahan menghadapi
        kembar identik yang berbagi pembagian namun berbeda?
\end{enumerate}
\end{method}

\section{Latihan}

\begin{exercise}[$\star$]\label{exo:g9:unique-individuals:1}
Sebutkan ketiga mekanisme pada teorema keunikannya.
\end{exercise}

\begin{exercise}[$\star$]\label{exo:g9:unique-individuals:2}
Apa yang ditambahkan pertukaran ruas selama pembagian duanya
pada pengocokannya?
\end{exercise}

\begin{exercise}[$\star$]\label{exo:g9:unique-individuals:3}
Nyatakan kedua paruh \cref{rem:g9:unique-individuals:alike}
--- persentasenya dan keunikannya --- dalam satu kalimat.
\end{exercise}

\begin{exercise}[$\star$]\label{exo:g9:unique-individuals:4}
Bagaimana kembar identik akhirnya tetap dapat dibedakan?
\end{exercise}

\begin{exercise}[$\star$]\label{exo:g9:unique-individuals:5}
Dua pertanyaan apa yang dijawab perbandingan \omterm{def:g9:genes-and-dna:dna}{DNA} bagi sebuah
pengadilan?
\end{exercise}

\begin{exercise}[$\star$]\label{exo:g9:unique-individuals:6}
Mengapa pendonor cangkok dicari dari keluarganya lebih dahulu?
\end{exercise}

\begin{exercise}[$\star\star$]\label{exo:g9:unique-individuals:7}
Kalikan hujahnya: gabungkan kedua ruang setengah tumpukan
orang tuanya dan lotere \omterm{def:g5:flower-to-fruit:steps}{pembuahannya} menjadi kesimpulan
``belum pernah dibagikan sebelumnya''.
\end{exercise}

\begin{exercise}[$\star\star$]\label{exo:g9:unique-individuals:8}
Uangkan \cref{prop:g9:unique-individuals:reserve} pada kebun
anggur \cref{exo:g8:sexual-reproduction:10}: apa yang tidak
dimiliki klonanya yang dimiliki \omterm{def:g8:reproduction-and-environment:population}{populasi} yang tumbuh dari \omterm{def:g2:from-seed-to-plant:seed}{biji}?
\end{exercise}

\begin{exercise}[$\star\star$]\label{exo:g9:unique-individuals:9}
Jalankan \cref{met:g9:unique-individuals:claims} pada:
``sebuah \omterm{def:g9:genes-and-dna:gene}{gen} bagi keberhasilan sekolah ditemukan''.
\end{exercise}

\begin{exercise}[$\star\star$]\label{exo:g9:unique-individuals:10}
Mengapa ``tumpukan manusia yang normal'' adalah frasa tanpa
acuan? Jawablah dengan cadangannya dan teoremanya.
\end{exercise}

\begin{exercise}[$\star\star$]\label{exo:g9:unique-individuals:11}
Sebuah uji keayahan membersihkan nama seorang ayah yang
diragukan dan dituduh keliru oleh tetangga pada
\cref{exo:g9:heredity-and-traits:11}. Jelaskan apa yang
diperiksa ujinya, kedudukan demi kedudukan.
\end{exercise}

\begin{exercise}[$\star\star\star$]\label{exo:g9:unique-individuals:12}
``Semuanya berbeda, semuanya sekerabat, dan keduanya oleh
mekanisme yang sama.'' Tuliskan alinea penutupnya:
pengocokan, \omterm{prop:g9:genes-and-dna:explains}{mutasi} dan sandi bersamanya, masing-masing
ditempatkan.
\end{exercise}

\section{Soal: Seminar Berkas Dingin}

\begin{problem}[{Soal akhir pekan --- satu helai rambut,
tiga pertanyaan, tiga puluh tahun}]\label{pb:g9:unique-individuals:1}
Sebuah seminar berkas dingin (rekaan): sebuah pembobolan
berumur tiga puluh tahun yang belum terpecahkan, satu helai
rambut yang terawetkan beserta \omterm{def:g5:first-look-at-cells:cell}{sel} \omterm{def:g1:plants-around-us:parts}{akarnya}, dan cara
perbandingan modern. Seminarnya menuntun muridnya menelusuri
apa yang dapat dan tidak dapat dikatakan pengenalan lewat \omterm{def:g9:genes-and-dna:dna}{DNA}.

\medskip\noindent\textbf{Bagian I --- Apa yang dimuat
rambutnya.}
\begin{enumerate}
  \item Mengapa \emph{\omterm{def:g1:plants-around-us:parts}{akar}} rambutnya harus terawetkan untuk
        sebuah pembacaan yang lengkap
        (\cref{def:g6:cell-unit-of-life:parts} tahu di mana
        \omterm{def:g9:genes-and-dna:dna}{DNA} tinggal)?
  \item Laboratoriumnya membaca seperangkat kedudukan yang
        sangat berubah-ubah. Mengapa yang berubah-ubah ---
        apa yang akan ditegakkan kedudukan yang 99 persen
        bersamanya?
  \item Profilnya cocok dengan seorang tersangka pada setiap
        kedudukan. Nyatakan apa yang ditegaskan kecocokannya,
        dan satu pengecualian sistematisnya
        (\cref{ex:g9:unique-individuals:twins}).
  \item Pembelanya mencatat bahwa tersangkanya punya seorang
        kembar identik di luar negeri. Apa yang kini harus
        dilakukan pengadilannya, dan mengapa?
\end{enumerate}

\medskip\noindent\textbf{Bagian II --- Pertanyaan
keluarganya.}
\begin{enumerate}[resume]
  \item Sebuah profil sebagian yang lebih awal menunjuk
        kepada ``seorang kerabat dekat keluarga K''.
        Jelaskan bagaimana kecocokan \emph{sebagian} terbaca
        sebagai \omterm{prop:g6:classification-kinship:kinship}{kekerabatan} (pembagian pada
        \cref{prop:g9:unique-individuals:theorem},
        dijalankan mundur).
  \item Seorang juri bertanya apakah profilnya menyingkapkan
        \omterm{def:g9:heredity-and-traits:traits}{sifat} tersangkanya --- rupanya, perangainya, ``\omterm{def:g9:genes-and-dna:gene}{gen}
        penjahatnya''. Jalankan
        \cref{met:g9:unique-individuals:claims} pada
        pertanyaan itu, pemeriksaan demi pemeriksaan.
  \item Juri yang lain bertanya apakah kecocokannya mungkin
        ``sekali di kota ini'' alih-alih sekali di seluruh
        umat manusia. Apa yang menetapkan jumlah kedudukan
        yang dibaca?
  \item Pemimpin seminarnya memperingatkan: ``\omterm{def:g1:living-or-not:living}{biologinya}
        berkata \omterm{def:g5:first-look-at-cells:cell}{sel} siapa; perkaranya masih harus berkata
        bagaimana \omterm{def:g5:first-look-at-cells:cell}{sel} itu sampai ke sana.'' Bedakan keduanya.
\end{enumerate}

\medskip\noindent\textbf{Bagian III --- Penutup
seminarnya.}
\begin{enumerate}[resume]
  \item Ringkaslah alatnya dalam satu kalimat: mekanisme dua
        bab yang mana (\cref{ch:g9:chromosomes},
        \cref{ch:g9:genes-and-dna}) yang membuat sebuah
        profil unik dan sebuah keluarga terbaca.
  \item Mengapa alat yang sama bekerja pada rambut berumur
        tiga puluh tahun? (\omterm{def:g9:heredity-and-traits:traits}{Sifat} \omterm{def:g9:genes-and-dna:dna}{DNA} yang mana --- di luar
        penyalinan pada \cref{prop:g9:genes-and-dna:explains}
        --- yang sedang disandari: kemantapannya.)
  \item Seminarnya berakhir pada etika: sebutkan dua
        pemakaian keunikan yang diperlihatkan bab ini
        melayani manusia, dan penyalahgunaan bersejarah yang
        diperingatkan \cref{rem:g9:unique-individuals:alike}.
  \item Tutuplah dengan teoremanya yang dinyatakan ulang
        untuk ruang sidangnya: apa yang tidak pernah terjadi
        dua kali, dan mengapa.
\end{enumerate}
\end{problem}
